% 20260602_Assingment_AppliedMathmatics


\title{20260602_応用解析学課題}
\date{}

\begin{document}
\begin{tcolorbox} %\begin{souquiz}
$A$を正定値行列とする。このとき$x= (x_1, \ldot, x_H)$とおくと。
\begin{equation*}
    I= \frac{1}{(2\pi)^\frac{H}{2}}\int_{-\infty}^\infty\cdot \int_{-\infty}^\infty \exp{-\frac{\tr{(x- \mu)}A^{-1}(x- \mu)}{2}}dx_1\cdots dx_H= 1
\end{equation*}
    であることを示す。
\end{tcolorbox} %\end{souquiz}

\begin{answer}
    $A$は正定値行列なので、直行行列$P$を用いて、
\begin{equation*}
    P^{-1}AP= 
\begin{pmatrix}
    \lambda_1& & \\
    & \vdots& \\
    & & \lambda_H
\end{pmatrix},\ \ (\lambda_i> 0)
\end{equation*}    
    とできる。ここで$P$か直行行列であることから$P^{-1}= \tr{P}$で$\det{P}= \det{P^{-1}}= \pm 1$を得るので、
\begin{equation*} \label{eq:detの定義}
    \det{A}= \prod_{i= 1}^H \lambda_i
\end{equation*}
    また、$y= \tr{P}(x- \mu)$とおけば、$y= (y_1, \ldots, y_H)$
\begin{equation*}
    \tr{(x- \mu)}A^{-1}(x- \mu)= \tr{y}\tr{P}A^{-1}Py
\end{equation*}
    ここで$(\tr{P}A^{-1}P)\cdot (P^{-1}AP)= \tr{P}\cdot P= I$より
\begin{align} \label{eq:行列をyで挟む}
    \tr{(x- \mu)}A^{-1}(x- \mu)
    &= \tr{y}
    \begin{pmatrix}
    1/\lambda_1& & \\
    & \vdots& \\
    & & 1/\lambda_H
    \end{pmatrix}
    y\\
    &= \sum_{i= 1}^H \frac{1}{\lambda_i}\cdot y^2_i
\end{align}
    を得られる。\zcref{eq:detの定義}、\zcref{eq:行列をyで挟む}を持ちいると、$y= \tr{P}(x- \mu)$のヤコビアンは$(\det{\tr{P}})$なので、
\begin{align}
    \frac{1}{(2\pi)^\frac{H}{2}\sqrt{\det{A}}}\int_{-\infty}^\infty\cdots \int_{-\infty}^\infty \exp{-\frac{\tr{(x- \mu)}A^{-1}(x- \mu)}{2}}dx_1\cdots dx_H
    &= \frac{1}{(2\pi)^\frac{H}{2}\sqrt{\prod_{i= 1}^H \lambda_i}}\int_{-\infty}^\infty\cdots \int_{-\infty}^\infty \exp{\sum_{i= 1}^H\left(frac{y_i}{2\lambda_i}\right)\left|\det{\tr{P}}\right|dy_1\cdots dy_H
\end{align}
    を得られる。ここで$|\det{\tr{P}}|= 1$と$z^2_i= \frac{y^2_i}{2\lambda_i}$とおけば、$\sqrt{2\lambda_i}dz_i= dy_i$なので各$\lambda_i> 0$より積分区間は交わらず、
\beign{align}
    I
    &= \frac{1}{(2\pi)^\frac{H}{2}\sqrt{\prod_{i= 1}^H \lambda_i}}\int_{-\infty}^\infty\cdots \int_{-\infty}^\infty \exp{\sum_{i= 1}^H\left(\sum_{i= 1}^H z^2_i\right)\cdot \prod_{i= 1}^H \sqrt{2\lambda_i}dz_1\cdots dz_H\\
    &= \frac{1}{\pi^\frac{H}{2}}\int_{-\infty}^\infty e^{z^2_H}\int_{-\infty}^\infty e^{z^2_{H- 1}}\cdots \int_{-\infty}^\infty e^{z^2_1}dz_1dz_2\cdots dz_{H- 1}dz_h
\end{align}
    と変形できる。Gauss積分$\int_{-\infty}^\infty e^{-x^2}dx= \sqrt{\pi}$であるから、
\begin{equation}
    I
    &= \frac{1}{\pi^\frac{H}{2}}\cdot \sqrt{\pi}\sqrt{\pi}\cdots \sqrt{\pi}\\
    &= \frac{\pi^\frac{H}{2}}{\pi^\frac{H}{2}}\\
    &= 1 
\end{equation}
    となり、
\begin{equation}
    \frac{1}{(2\pi)^\frac{H}{2}}\int_{-\infty}^\infty\cdot \int_{-\infty}^\infty \exp{-\frac{\tr{(x- \mu)}A^{-1}(x- \mu)}{2}}dx_1\cdots dx_H= 1
\end{equation}
    が示された。
\end{answer}
\end{document}
